\section{Fugitive emissions}
\label{sec:fugitive}

\subsection{Scope and inventory mapping}

This block targets the \textbf{area-source} component of CAMS GNFR~D
(\emph{Fugitive emissions}). In the CAMS-REG-ANT inventory, GNFR~D is
reported as one aggregated sector, even though the underlying emission
processes include solid-fuel extraction, oil transport and upstream
operations, storage and refining activities, gas handling, flaring, and
other residual energy-related fugitive sources. The methodological
problem at this stage is therefore not the construction of final
downscaled emissions, but the construction of pollutant-specific
\emph{within-cell spatial weights} that can later be combined with the
corresponding CAMS cell totals.

This distinction is essential. The workflow described here produces a
multi-band raster of normalised fine-grid weights, one band per
pollutant, such that the fine-pixel shares sum to unity within each
CAMS GNFR~D area cell. It does not by itself represent emissions, nor
does it assume that CAMS already resolves the internal fugitive
subsector structure required for spatial allocation.

\subsection{Conceptual design}

GNFR~D is internally heterogeneous. Coal mining, oil terminals and
pipelines, refinery and storage sites, and gas infrastructure do not
share the same fine-scale geography. A single proxy would therefore
force all fugitive-emission pollutants to follow the same spatial
pattern despite their different process origins. That simplification
would introduce strong spatial bias, especially where one country is
dominated by one fugitive process family while another is dominated by a
different one.

The upgraded design therefore separates the problem into three linked
elements:
\begin{itemize}
  \item a set of four operational proxy groups representing the main
        spatial environments of fugitive emissions;
  \item pollutant-specific country-level group shares derived from
        CEIP/NRD emissions;
  \item a final composite proxy that recombines the common group
        surfaces with the pollutant-specific national shares and is then
        normalised within each CAMS parent cell.
\end{itemize}

This produces a transparent two-stage weighting system. First, the
ancillary geospatial datasets describe where different fugitive
activities are most plausibly located. Second, the CEIP/NRD national
totals determine how strongly each pollutant should follow each group in
each country. Pollutant bands may therefore differ spatially even though
CAMS provides only one aggregate GNFR~D area layer.

\subsection{Operational groups and pollutant shares}

For the purpose of spatial weighting, the detailed CEIP/NRD fugitive
codes are aggregated into four operational groups:
\begin{table}[htbp]
  \centering
  \caption{Operational groups used for the GNFR~D fugitive weighting workflow.}
  \label{tab:fugitive_groups}
  \begin{tabular}{p{1.2cm}p{3.2cm}p{7.0cm}}
    \hline
    Group & CEIP/NRD codes & Intended spatial interpretation \\
    \hline
    G1 & 1B1a, 1B1b, 1B1c & Coal and solid-fuel extraction environments \\
    G2 & 1B2ai & Oil transport and upstream infrastructure \\
    G3 & 1B2aiv, 1B2av & Storage, refining, and fuel distribution environments \\
    G4 & 1B2b, 1B2c, 1B2d & Gas handling, flaring, and residual fugitive-energy environments \\
    \hline
  \end{tabular}
\end{table}

The operational split is intentionally spatial rather than chemical. Its
purpose is to identify groups with distinct geographies that can be
proxied separately and then recombined. In the implemented grouping, all
of CEIP code 1B2b is assigned to G4.

Let $c$ denote country and $p$ pollutant. The national CEIP/NRD group
totals are written as
\begin{equation}
  E_{G1,c,p},\qquad E_{G2,c,p},\qquad E_{G3,c,p},\qquad E_{G4,c,p}.
\end{equation}
The corresponding pollutant-specific group shares are
\begin{equation}
  \alpha_{g,c,p} =
  \frac{E_{g,c,p}}
       {\displaystyle\sum_{g' \in \{G1,G2,G3,G4\}} E_{g',c,p}}.
\end{equation}
By construction,
\begin{equation}
  \sum_{g \in \{G1,G2,G3,G4\}} \alpha_{g,c,p} = 1.
\end{equation}

Whenever pollutant-specific group totals are available, the
$\alpha_{g,c,p}$ coefficients are computed directly for each
country--pollutant pair. If a pollutant has no usable group split for a
given country, the implementation falls back to total fugitive
emissions across pollutants for that same country and derives a common
group mixture from those totals. The fallback therefore preserves the
country-specific fugitive structure whenever pollutant-specific detail is
missing.

\subsection{Group-specific spatial proxies}

Let $i$ denote a fine pixel on the reference grid and let $A_i$ be its
area. The method constructs four non-negative group surfaces
$P_{G1}(i),\ldots,P_{G4}(i)$ from harmonised ancillary datasets. These
surfaces do not represent emissions directly; they represent the
relative spatial plausibility of each fugitive group.

\subsubsection{General formulation}

Each spatial component is first converted to a unitless relevance field
using a clipped quantile scaling:
\begin{equation}
  z(x) =
  \mathrm{clip}_{[0,1]}
  \left(
    \frac{
      \min\!\left(\max(x, q^{0.01}), q^{0.99}\right) - q^{0.01}
    }{
      q^{0.99} - q^{0.01}
    }
  \right),
\end{equation}
where $q^{0.01}$ and $q^{0.99}$ are the 1st and 99th percentiles of the
finite values of the corresponding layer. Missing values are set to
zero. This scaling preserves relative spatial contrasts while preventing
a small number of extreme pixels from dominating the proxy.

The population component is defined from population density:
\begin{equation}
  P_{\mathrm{pop}}(i) =
  z\!\left(\frac{\mathrm{Pop}(i)}{A_i}\right).
\end{equation}
It acts both as a diffuse-support term and as a fallback proxy where the
sector-specific layers are absent.

For each group $g$, the OSM component is defined as the fractional area
share of selected fugitive polygons within the fine cell:
\begin{equation}
  OSM_g(i) =
  \frac{
    \mathrm{area}\!\left(
      \text{selected OSM polygons for group } g \cap i
    \right)
  }{
    \mathrm{area}(i)
  }.
\end{equation}
This term captures high-confidence infrastructure geometry such as
extraction sites, terminals, refinery polygons, storage areas, and gas
facilities.

The CORINE component provides broader eligibility and background
continuity. If $\mathcal{C}_g$ denotes the set of CORINE classes allowed
for group $g$, the class-support field is written as
\begin{equation}
  CLC_g(i) =
  \frac{1}{|\mathcal{C}_g|}
  \sum_{\ell \in \mathcal{C}_g}
  \mathbf{1}\!\left(\mathrm{CLC}(i)=\ell\right).
\end{equation}
Because CORINE classes are mutually exclusive at pixel scale, this term
acts as a blended eligibility score for multi-class groups rather than
as a true sub-pixel land-cover fraction.

The sector-specific group proxy is then
\begin{equation}
  P_{\mathrm{sector},g}(i) =
  0.6 \cdot z\!\left(OSM_g(i)\right)
  + 0.4 \cdot z\!\left(CLC_g(i)\right),
\end{equation}
and the full group proxy is
\begin{equation}
  P_g(i) =
  0.8 \cdot P_{\mathrm{sector},g}(i)
  + 0.2 \cdot P_{\mathrm{pop}}(i).
\end{equation}

This formulation assigns distinct roles to the three ingredients. OSM
captures high-confidence fugitive infrastructure, CORINE ensures spatial
continuity beyond the mapped polygons themselves, and population
provides a diffuse support term for emissions not directly resolved by
the fugitive geospatial layers.

\subsubsection{Group definitions}

\paragraph{G1 -- Coal and solid fuels.}
This group targets coal-mining and related solid-fuel environments. The
CORINE support uses class~131 (mineral extraction sites) and class~121
(industrial or commercial units). The OSM filter retains mines,
quarries, coal-related extraction areas, and other industrial polygons
consistent with solid-fuel fugitive activity.

\paragraph{G2 -- Oil transport and upstream operations.}
This group targets oil terminals, pipeline-related facilities, and
upstream oil infrastructure. The CORINE support uses class~123 (port
areas) and class~121 (industrial or commercial units). The OSM filter
retains oil infrastructure, terminals, pipeline-related polygons, and
port-related energy facilities.

\paragraph{G3 -- Storage, refining, and distribution.}
This group represents the storage and refining environment of fugitive
oil-product handling. The CORINE support uses class~121 (industrial or
commercial units) and class~122 (road and rail networks and associated
land). The OSM filter retains refineries, storage tanks, depots, fuel
terminals, and related distribution infrastructure.

\paragraph{G4 -- Gas handling, flaring, and residual fugitive energy.}
This group targets gas infrastructure, flaring, and residual
energy-related fugitive processes. The CORINE support uses class~121
(industrial or commercial units) and class~131 (mineral extraction
sites). The OSM filter retains flaring locations, gas facilities, power
plants, and related energy infrastructure. In the operational grouping
used here, code 1B2b is fully assigned to this group.

\subsection{Fallback logic}

If the sector-specific support for group $g$ is absent in a given fine
pixel, the method falls back completely to the population proxy. In
practice, if
\begin{equation}
  OSM_g(i) + CLC_g(i) = 0,
\end{equation}
then
\begin{equation}
  P_g(i) = P_{\mathrm{pop}}(i).
\end{equation}
This is a full replacement rather than a partial mixture. The fallback
ensures that CAMS parent cells retain a complete spatial support even
when neither mapped fugitive polygons nor eligible CORINE classes are
present locally.

\subsection{Composite proxy and within-cell normalisation}

Let $m$ denote a CAMS GNFR~D area cell and let $i \in m$ denote fine
pixels whose centres fall inside that parent cell. Each group surface is
first normalised within the CAMS cell:
\begin{equation}
  W_g(i) =
  \frac{P_g(i)}
       {\displaystyle\sum_{k \in m} P_g(k)}.
\end{equation}
If the denominator is zero for a given group and CAMS cell, the
implementation falls back to a uniform distribution over the valid
fine-grid support of that CAMS cell.

The pollutant-specific composite fugitive proxy is then
\begin{equation}
  P_{D,c,p}(i) =
  \sum_{g \in \{G1,G2,G3,G4\}}
  \alpha_{g,c,p} \cdot W_g(i),
\end{equation}
where $c$ denotes the country associated with pixel $i$.

Because the composite combines several already normalised group fields,
it is normalised a second time within each CAMS parent cell:
\begin{equation}
  W_{D,c,p}(i) =
  \frac{P_{D,c,p}(i)}
       {\displaystyle\sum_{k \in m} P_{D,c,p}(k)}.
\end{equation}
By construction,
\begin{equation}
  \sum_{i \in m} W_{D,c,p}(i) = 1
  \qquad \forall m,p.
\end{equation}

The final product is therefore a set of pollutant-specific
\emph{normalised weighting rasters}. These layers are intended for a
later downscaling step in which they may be combined with the
corresponding CAMS cell totals. They are not themselves the downscaled
emissions.

\subsection{Output and interpretation}

The output of the workflow is a fine-grid spatial weighting raster
aligned with the common proxy grid used by the other EL sectors. One
layer is written for each enabled CAMS pollutant. Each layer contains
relative within-cell weights only: it identifies how the CAMS GNFR~D
area-source total should be distributed inside each parent CAMS cell,
but it does not alter the total mass reported by CAMS.

The interpretation of the fugitive weighting raster is therefore
strictly allocative. OSM polygons, CORINE classes, and population
density are used only to describe plausible spatial variation in the
location of fugitive sources. The raster should not be read as a direct
estimate of extraction throughput, fuel handling, or emissions in
physical units at pixel scale.

\subsection{Limitations}

The GNFR~D workflow remains a first-order allocation scheme for a
heterogeneous sector. Several limitations follow directly from that
role.

First, CAMS provides only an aggregate fugitive area layer and does not
resolve the internal operational groups used here. The four-group split
is therefore an external analytical construction designed to recover more
credible spatial diversity from the aggregated inventory.

Second, the OSM fugitive layer is a major improvement over a pure
land-cover proxy, but its completeness and tagging quality are not
uniform. Missing polygons or inconsistent tags may suppress group
signals in some locations even where fugitive infrastructure is present.

Third, CORINE provides necessary spatial continuity, but it remains a
background land-cover product rather than a direct observation of
fugitive infrastructure or activity intensity. Its role is therefore to
stabilise the proxy, not to recover subsector-specific infrastructure on
its own.

Fourth, the use of population as a fallback is methodologically useful
because it prevents zero-support cells, but it should be interpreted as
a diffuse regularisation device rather than as evidence that fugitive
emissions follow settlement density directly.

Despite these limitations, the GNFR~D workflow represents a substantial
improvement over a single undifferentiated fugitive proxy. It preserves
country- and pollutant-specific group structure from CEIP/NRD reporting,
allows coal, oil, storage, and gas-related fugitive processes to follow
different fine-scale geographies, and maintains exact within-CAMS-cell
normalisation for the weighting raster used later in the downscaling
pipeline.
